\section{Conclusion}
\label{sec:conclusion}

Three overarching conclusions emerge from this empirical audit regarding physics-constrained wind power forecasting.

First, at the literature and reporting level, studies in the audited scoping cohort reported exclusively within nominal operating envelopes (0/46 documenting failure regimes), while frequently relying on confounded module ablations and omitting public code. When computational graphs are subjected to static and gradient verification, implementations can diverge substantially from published equations. Establishing reliable scientific claims in physics-constrained machine learning requires open source code, computational graph verification, and form-matched negative controls.

Second, at the methodological and inductive level, physical priors function not as universal boosters, but as \textbf{conditional inductive guardrails}. In restricted training windows, authentic aerodynamic curves supply vital inductive regularization that curtails tail dispatch risk, exhibiting prior-specific mapping effects in this sample-restricted setting that cannot be attributed to data volume alone. As operational telemetry expands, purely data-driven models match or slightly exceed fixed priors---a dynamic interpreted through the lens of diminishing variance reduction and emerging structural approximation bias, though these mechanisms remain conceptual interpretations consistent with---not identified by---the observed learning curves. Furthermore, physical priors exhibit pronounced fragility under uncalibrated engineering drift and stale telemetry proxies (where decision risk exceeds the unconstrained reference by the first evaluated lag $k=6$ / 1\,h delay and deteriorates monotonically through $k=36$). Operational deployment therefore mandates adaptive regularization annealing and telemetry latency tolerance checks.

Third, at the operational decision level, symmetric regression metrics (RMSE and MAE) and standalone unhedged balancing cashflows structurally decouple from real-world dispatch risk. Symmetric metrics mask catastrophic tail under-predictions during shortfall events, while unhedged single-component settlement rewards degenerate, non-forecasting strategies. Forecasting models intended for grid integration must be optimized and audited against decision-aligned asymmetric dispatch objectives.

In conclusion, physics-constrained machine learning for energy systems must move beyond uncritical performance narratives. By institutionalizing form-matched negative controls, boundary stress tests, computational graph auditing, and decision-aligned evaluation (Criteria C1--C5), the field can transition toward verifiable, reproducible, and operationally dependable forecasting systems.
